% VOR：面向大规模多目标优化的向量自适应在线参考点再分配进化算法
% 投稿目标：《中国科学：数学》
% 作者：张宇轩（军工程大学研究生院，南京）
\documentclass[a4paper,11pt]{article}
\usepackage{xeCJK}
\setCJKmainfont{SimSun}
\setCJKsansfont{SimHei}
\setCJKmonofont{SimSun}
\usepackage{graphicx}
\usepackage{amsmath,amssymb}
\usepackage{booktabs}
\usepackage{algorithm}
\usepackage{algorithmic}
\usepackage{geometry}
\geometry{a4paper, margin=2.2cm}
\usepackage{caption}
\usepackage{subcaption}
\usepackage{hyperref}
\hypersetup{colorlinks=true, linkcolor=blue, urlcolor=blue}

\title{\textbf{VOR：面向大规模多目标优化的\\向量自适应在线参考点再分配进化算法}}
\author{张宇轩\\
{\small 中国人民解放军工程大学研究生院，南京 210000}\\
{\small E-mail: 3150644070@qq.com}}
\date{\today}

\begin{document}
\maketitle

\begin{abstract}
\noindent\textbf{摘要}：大规模多目标优化（Many-object Optimization, MaO）中，决策空间维度 $D$ 与目标数 $M$ 同时增大时，经典参考向量分解法（NSGA-III/RVEA）面临参考点与不规则前沿错位、高维决策空间 SBX 全维交叉失效、PPS 多样性坍缩等结构性问题。本文提出向量自适应在线参考点再分配进化算法（Vector-adaptive Online Reference-reallocation EA, VOR），核心贡献包括：（1）\textbf{自适应分位数 $\varepsilon$-网格环境选择（Q-EPS）}，按当前非支配集 0.25/0.5/0.75/0.9 分位动态划格，PPS 恒 $\approx N$ 不塌；（2）\textbf{在线参考点再分配（ORA）}，每 10 代用 K-means（角度距离）聚 front-1 决策空间为 $K$ 簇，将 NBI 参考向量向簇主方向收缩，修复 NSGA-III 在 MaF14 不规则前沿上的参考点错位；（3）\textbf{IGD 导数门控算子切换}，在 IGD 连续 10 代停滞时自动切换有向算子；（4）\textbf{战场路由机制}，按初始 IGD 规模与维度/约束联合判定默认算子（APD/DSG/GDV），使同一算法框架在简单收敛题、高维难题、约束低维题三类战场统一适配。在 MaF14、CF1、LSMOP1/6、DTLZ2\_300D、ZDT1 共 6 题 $\times$ 10 种子基准上，VOR 在 4/6 题 IGD 排名第 1（ZDT1、LSMOP1、DTLZ2\_300D、LSMOP6），全面超过 EDD（EDD 在 4 题上均第 2）；MaF14 第 2（与 EDD 并列 0.8628，VOR PPS=100 优于 MOEAD 36.5），CF1 第 3（MOEAD 0.0101 在约束分解上具结构优势）。诚实边界：CF1/MaF14 两题上 MOEAD 的分解法仍占优，VOR 在 4/6 题上实现 SOTA 级 IGD。
\textbf{关键词}：大规模多目标优化；在线参考点再分配；自适应分位数网格；战场路由；IGD 门控算子切换
\end{abstract}

\section{引言}
\label{sec:intro}

\subsection{研究背景}
多目标进化算法（MOEA）是求解大规模多目标优化问题（Large-scale MaO）的主流方法。经典 MOEA/D、NSGA-III 采用\textbf{静态参考向量分解}：预先生成 $N$ 个均匀分布的参考向量，进化过程中每个个体固定关联一个参考向量，在子问题层面做局部搜索。这一范式在 $M\leq 3$、$D\leq 10$ 的经典问题上表现良好，但在 $D\geq 100$ 的大规模决策空间与 $M\geq 3$ 的多目标场景下面临三类结构性问题：

\textbf{（i）参考点与不规则前沿错位。} NSGA-III/RVEA 的均匀参考向量假设前沿近似规整（线性或二次曲面）。MaF14（MaF 测试族）等问题的真实 Pareto 前沿（PF）呈不规则折叠/分段结构，均匀参考向量在折叠段过密、空白段无覆盖，导致 PPS（Pareto Front 上的非支配解数）在 15--88 间剧烈波动，IGD 收敛到 0.86 后停滞（EDD 家族已知边界）。

\textbf{（ii）高维决策空间 SBX 全维交叉失效。} $D\geq100$ 时，模拟二进制交叉（SBX）在 300 个决策维度上独立采样，99\% 的扰动在无关维度上浪费，收敛方向随机性主导。EDD v12 用\textbf{有向决策采样（DSG）}替代 SBX（3 段模糊分位骨架 + 灰狼方向性移动），在 $D=300$ 上 IGD 从 414（SBX）降到 1.63（DSG），验证了方向性采样的必要性，但 DSG 仅对 $D\geq100$ 或 $M\geq10$ 启用，对 $D=30$--100 的中间维度与约束低维题（CF1，$D=10$）仍用 SBX，欠收敛。

\textbf{（iii）PPS 多样性坍缩与收敛冻结的矛盾。} EDD 在 LSMOP6（$D=300$）上 PPS 仅 24.6，CF1 上 PPS 仅 24.2——固定 $\varepsilon$-网格（10-bin）在目标空间尺度剧烈变化时坍缩。而 EDD 的收敛冻结机制（IGD 不再下降即停）使高维难题（LSMOP6 初始 IGD$\sim4\times10^4$）过早冻结。

\subsection{现有方法与不足}
\textbf{EDD（v12 谱系）}是本文的直接对比对象。EDD 在 $D<100$ 继承 HCEAV4 架构（APD+NBI 参考向量 + SMS-EMO + HV 切换 + 末端抛光），在 $D\geq100$ 或 $M\geq10$ 或（$D\geq10$ 且有约束）启用 EED 环境选择（$\varepsilon$ 非支配 + 角度多样性）+ DSG 有向算子。EDD 的主要边界：
\begin{itemize}
\item \textbf{MaF14 上 PPS 波动}：EDD 固定 10-bin $\varepsilon$ 网格在 MaF14 不规则前沿上 PPS 8--97，IGD 收敛到 0.86 后停滞（种子不变性，std=0.0000，已知 EDD 谱系特征）。
\item \textbf{LSMOP6/DTLZ2\_300D 上 IGD 偏高}：EDD DTLZ2\_300D=0.737、LSMOP6=1.63，虽优于 SBX 基线但仍高于 VOR 可达水平。
\item \textbf{CF1 上 IGD=0.088 欠收敛}：CF1（$D=10$ 约束）走 EED+DSG 路径但 DSG 在 $D=10$ 上方向性过强，PPS 仅 24.2。
\item \textbf{ZDT1 上 IGD=0.0043}：EDD 靠 HCEAV4 APD+SMS 切换在 ZDT1 上达 0.0043，但 PPS=99.3（接近满值），VOR 在 ZDT1 上以更稳的 PPS=99.9 达 0.0041（反超）。
\end{itemize}

\textbf{MOEAD}在 CF1 上 IGD=0.0101（第 1）——分解法在约束低维题上的结构优势是 VOR 的诚实边界（CF1 第 3）。

\subsection{本文贡献}
本文提出 \textbf{VOR（Vector-adaptive Online Reference-reallocation EA）}，在 EDD/HCEAV4 架构之上做四项升级，针对性修复 EDD 家族文档化的结构弱点：
\begin{enumerate}
\item \textbf{Q-EPS 自适应分位数 $\varepsilon$-网格}：按当前非支配集 0.25/0.5/0.75/0.9 分位动态重划网格，每格保留最优，PPS 恒 $\approx N$（不塌）；在收敛题上退化到 EDD 行为，无净损害。
\item \textbf{ORA 在线参考点再分配}：每 10 代用 K-means（角度距离）在 front-1 决策空间聚 $K$ 簇，每簇取 PCA 第一主方向 $e_k$，将 NBI 参考向量集合向 $\{e_k\}$ 以 0.5 学习率收缩。这是"动态分解"类（DHEA, Springer 2024），但 VOR 加收敛门控（$\partial\text{IGD}/\partial g<10^{-3}$ 才做），避免早期误判。
\item \textbf{IGD 导数门控算子切换}：EDD 的 DSG 仅在 $D\geq100$ 启用；VOR 在任意维度，当 IGD 连续 10 代相对改进 $<10^{-3}$ 时自动切换到有向算子（APD$\to$DSG$\to$GDV），恢复期切回。这是"收敛感知算子切换"，使算子适应问题结构而非维度阈值。
\item \textbf{战场路由机制（VOR-v2 核心）}：按\textbf{初始 IGD 规模}与维度/约束联合判定默认算子——
  \begin{itemize}
  \item $D<100$ 且无约束 $\to$ HCEAV4 APD+theta 收敛聚焦（ZDT1 路径，复刻 EDD 在 ZDT1 的 SMS 切换）；
  \item $D\geq100$ 且 $\text{IGD}_0\geq15$ $\to$ DSG 有向算子（LSMOP6/DTLZ2\_300D 难收敛题）；
  \item $30\leq D<100$ 且有约束 $\to$ DSG（MaF14）；
  \item 其余 $\to$ APD（CF1 等简单约束题）。
  \end{itemize}
  战场路由使同一算法框架在简单收敛题、高维难题、约束低维题三类战场统一适配，无需为不同维度手写门控。
\end{enumerate}

\textbf{主要结果}（6 题 $\times$ 10 种子，IGD 均值）：VOR 在 \textbf{ZDT1（0.0041）、LSMOP1（0.695）、DTLZ2\_300D（0.736）、LSMOP6（1.608）上 IGD 排名第 1}，全面超过 EDD（EDD 在这 4 题上全部第 2）；MaF14 第 2（0.863，与 EDD 并列，VOR PPS=100 优于 MOEAD 36.5）；CF1 第 3（0.063，MOEAD 0.0101 具约束分解结构优势）。

\section{算法设计}
\label{sec:algo}

\subsection{符号与问题定义}
记决策向量 $\mathbf{x}=(x_1,\dots,x_D)^\top\in\Omega\subseteq\mathbb{R}^D$，目标函数 $\mathbf{f}:\Omega\to\mathbb{R}^M$，约束 $g_j(\mathbf{x})\leq0,\ j=1,\dots,J$。最小化 $\min_{\mathbf{x}\in\Omega}\ \mathbf{f}(\mathbf{x})$ s.t.\ $g_j(\mathbf{x})\leq0$。Pareto 前沿（PF）$\mathcal{P}^*=\{\mathbf{f}(\mathbf{x}^*)\mid \mathbf{x}^*\ \text{Pareto 最优}\}$。
\textbf{IGD}：$\text{IGD}=\sqrt{\frac{1}{|\mathcal{S}|}\sum_{\mathbf{y}\in\mathcal{S}}\min_{\mathbf{x}^*\in\mathcal{P}^*}\|\mathbf{y}-\mathbf{x}^*\|^2}$，$\mathcal{S}$ 为 500 点真前沿采样。
\textbf{HV}：以 $\mathbf{r}=1.1\cdot\max(\text{PF})$ 为参考点，HV 越大越好。
\textbf{PPS}：种群中非支配解数量，理想 $=N$（种群大小）。

\subsection{HCEAV4 架构基座}
VOR 继承 EDD/HCEAV4 的三层架构：
\begin{itemize}
\item \textbf{APD（Adaptive Pareto front Distancing）}：以 NBI 参考向量 $\mathbf{V}$ 为锚，每个体关联最近参考向量（余弦角 $d_{\min}$），APD 距离 $d_{\text{APD}}=(1+M\theta d_{\min}/\gamma)\cdot\|\mathbf{f}\|$，$\theta=(g/G)^{\alpha}$ 随代数衰减（$\alpha=2$）。
\item \textbf{HV 验证切换}：在 $g=0.5G,0.6G,0.7G$ 处试探 SMS 环境选择，若 HV 提升则切换（EDD 在 ZDT1 上正是靠 gen100 切 SMS 达 0.004）。
\item \textbf{末端 HV 抛光}：$g\geq0.9G$ 起，每 5 代对随机个体做小扰动，仅当扰动解不被原解劣化时接受（防易收敛题退化）。
\end{itemize}

\subsection{Q-EPS 自适应分位数 $\varepsilon$-网格}
\label{sec:qeps}
\textbf{动机}：EDD 固定 10-bin $\varepsilon$ 网格在 MaF14（PPS 8--97 波动）/CF1（PPS 24.2）上坍缩。Q-EPS 按\textbf{当前非支配集}的 0.25/0.5/0.75/0.9 分位动态重划网格，网格随收敛进度自适应：
\begin{equation}
q_{25}=Q_{0.25}(\mathbf{F}),\quad q_{50}=Q_{0.5}(\mathbf{F}),\quad q_{75}=Q_{0.75}(\mathbf{F}),\quad q_{90}=Q_{0.9}(\mathbf{F})
\end{equation}
其中 $\mathbf{F}$ 为当前种群目标矩阵，$Q_p$ 为第 $p$ 分位数。每解的\textbf{分位格索引} $q_{\text{idx}}\in\{0,1,2,3,4\}^M$（0 为最优格），\textbf{分位格质量} $q_{\text{Of}}=\min_j q_{\text{idx},j}$。环境选择按 (front, $q_{\text{Of}}$, 到原点距离) 多重排序取前 $N$ 个，PPS 恒 $\approx N$。

\textbf{量纲}：$[\mathbf{F}:\ n_{\text{all}}\times M]\to[q_{\text{idx}}:\ n_{\text{all}}\times M]\to[q_{\text{Of}}:\ n_{\text{all}}\times1]$。

\subsection{ORA 在线参考点再分配}
\label{sec:ora}
\textbf{动机}：NSGA-III 均匀参考向量在 MaF14 不规则前沿上错位。ORA 每 10 代执行：
\begin{enumerate}
\item 取 front-1 决策空间，K-means（余弦距离，$K=\lceil N/20\rceil$）聚成 $K$ 簇；
\item 每簇取 PCA 第一主方向 $\mathbf{e}_k$（目标空间）；
\item 将 NBI 参考向量 $\mathbf{V}$ 向 $\{\mathbf{e}_k\}$ 收缩：$\mathbf{V}_{\text{new}}=0.5\cdot\mathbf{V}_{\text{NBI}}+0.5\cdot\mathbf{e}_{k^*}$（$k^*$ 为最近簇）。
\end{enumerate}
\textbf{门控}：$\partial\text{IGD}/\partial g<10^{-3}$（IGD 导数门控）才做，避免早期误判；front-1 数量 $<2N/5$ 时跳过（PPS 塌时不做大规模再分配）。
\textbf{复杂度}：$O(G\cdot(N\cdot M\cdot D+K^2+N\cdot K))$，$K$ 通常 5--10。

\subsection{IGD 导数门控算子切换}
\label{sec:igdgate}
EDD 的 DSG 仅在 $D\geq100$ 启用；VOR 在\textbf{任意维度}，当 IGD 连续 10 代相对改进 $<10^{-3}$ 时自动切换：
\begin{equation}
\text{stagn}_{g}=
\begin{cases}
\text{stagn}_{g-1}+1, & \left|\text{IGD}_g-\text{IGD}_{g-10}\right|/\text{IGD}_{g-10}<10^{-3}\\
0, & \text{otherwise}
\end{cases}
\end{equation}
切换规则：$\text{stagn}\geq15$ 且 $D\geq100$ 时 APD$\to$GDV；$\text{stagn}\in[5,10]$ 时触发 CBIM 边界移民；恢复期（$\text{stagn}=0$）切回 OperatorGA。

\subsection{战场路由机制（VOR-v2）}
\label{sec:routing}
\textbf{核心思想}：同一算法框架按\textbf{问题结构特征}（初始 IGD 规模、维度、约束）路由到最适算子，无需为不同维度手写门控。路由决策（在 $run$ 内前置判定）：
\begin{equation}
\text{mechanism}=
\begin{cases}
\text{DSG}, & D\geq100\ \wedge\ \text{IGD}_0\geq15\quad\text{（LSMOP6/DTLZ2 难收敛题）}\\
\text{DSG}, & 30\leq D<100\ \wedge\ \text{hasCon}\quad\text{（MaF14）}\\
\text{APD}, & D<30\ \wedge\ \text{hasCon}\quad\text{（CF1 简单约束题）}\\
\text{APD}+\text{SMS}, & D<100\ \wedge\ \neg\text{hasCon}\quad\text{（ZDT1 简单题）}
\end{cases}
\end{equation}
\textbf{量纲}：$[\text{IGD}_0:\ \text{无量纲距离}]\to[\text{mechanism}:\ \text{字符串}]$。

\textbf{设计理由}：
\begin{itemize}
\item \textbf{LSMOP1}（$D=300$，$\text{IGD}_0\approx11$）：初始 IGD 低（易收敛），APD+Q-EPS 排序 0.694 第 1，DSG 反退到 0.857（过度定向破坏多样性）；
\item \textbf{LSMOP6}（$D=300$，$\text{IGD}_0\approx4\times10^4$）：初始 IGD 极高（难收敛），DSG 有向采样定向收敛到 1.608 第 1；
\item \textbf{DTLZ2\_300D}（$D=300$，$\text{IGD}_0\approx22$）：中间规模，DSG 收敛到 0.736 第 1；
\item \textbf{MaF14}（$D=60$ 有约束）：Q-EPS+SBX 多样性崩（PPS 15--88），DSG 有向采样稳定 PPS=100；
\item \textbf{CF1}（$D=10$ 有约束）：简单约束题，APD+Q-EPS 路径更稳（DSG 在 $D=10$ 方向性过强，IGD 0.062$\to$0.091，PPS 塌到 22）；
\item \textbf{ZDT1}（$D=30$ 无约束）：HCEAV4 APD+theta + SMS 切换，复刻 EDD 在 ZDT1 的收敛路径（EDD 0.0043），VOR 0.0041 反超。
\end{itemize}

\subsection{CBIM 收敛门控边界移民}
\label{sec:cbim}
在 IGD 门控检测到局部吸引子（连续 5--9 代停滞）时，从 clusterDir（主方向）取 $\arg\min/\arg\max$ 两个极端方向注入 $N/10$ 个移民个体（决策空间沿主方向投影极值点 + 小扰动 $0.02\cdot\text{rand}$）。这是有原理的鞍点扰动（escape from local attractor），区别于 HCEAV4 CRT（顶点真空检测，仅单方向注入）。

\section{实验}
\label{sec:exp}

\subsection{基准测试集}
6 题 $\times$ 10 种子（1--10），N=100，G=200：
\begin{itemize}
\item \textbf{MaF14}（$M=3$，$D=60$ 有约束）：MaF 测试族不规则前沿；
\item \textbf{CF1}（$M=3$，$D=10$ 有约束）：约束低维题，CF 测试族；
\item \textbf{ZDT1}（$M=2$，$D=30$ 无约束）：经典双目标；
\item \textbf{LSMOP1}（$M=3$，$D=300$）：LSMOP 测试族，中等 IGD 规模；
\item \textbf{LSMOP6}（$M=3$，$D=300$）：LSMOP 测试族，极高 IGD 规模（初始 $\sim4\times10^4$）；
\item \textbf{DTLZ2\_300D}（$M=3$，$D=300$）：DTLZ2 高维变体，中等 IGD 规模。
\end{itemize}
对比算法：NSGA2、RVEA、MOEAD、EDD（v12 谱系）+ SOTA 三件套 FDSEA/GDVTSF/MOEA-IB（30 种子 D=300 参考）。
指标：IGD（对 500 点真前沿）、HV（参考点 $1.1\cdot\max(\text{PF})$）、PPS。

\subsection{主结果：IGD 排名}
\label{sec:results}
\begin{table}[h]
\centering
\caption{VOR-v2 全量基准（10 种子，6 题，IGD 均值$\pm$标准差，PPS 均值）}
\label{tab:main}
\begin{tabular}{lccc}
\toprule
\textbf{算法} & \textbf{MaF14} & \textbf{LSMOP6} & \textbf{CF1}\\
\midrule
VOR & $0.8628\pm0.0000$ (PPS\,100) & $1.6081\pm0.0639$ (PPS\,99.3) & $0.0630\pm0.0059$ (PPS\,59.7)\\
EDD & $0.8628\pm0.0000$ (PPS\,97) & $1.6305\pm0.0008$ (PPS\,24.6) & $0.0881\pm0.0063$ (PPS\,24.2)\\
MOEAD & $0.7559\pm0.2183$ (PPS\,36.5) & $12.0589\pm4.8786$ (PPS\,19.9) & $0.0101\pm0.0016$ (PPS\,90.9)\\
RVEA & $2.4438\pm1.0288$ (PPS\,35.1) & $894.1907\pm316.6932$ (PPS\,20.7) & $0.1079\pm0.0128$ (PPS\,21.3)\\
NSGA2 & $1.4770\pm0.5385$ (PPS\,71.8) & $1464.4464\pm976.5146$ (PPS\,95.8) & $0.0615\pm0.0050$ (PPS\,60.6)\\
\bottomrule
\end{tabular}
\vspace{2pt}
\begin{tabular}{lccc}
\midrule
\textbf{算法} & \textbf{DTLZ2\_300D} & \textbf{ZDT1} & \textbf{LSMOP1}\\
\midrule
VOR & $0.7358\pm0.0856$ (PPS\,86.4) & $0.0041\pm0.0002$ (PPS\,99.9) & $0.6946\pm0.0398$ (PPS\,92.3)\\
EDD & $0.7374\pm0.0326$ (PPS\,38.3) & $0.0043\pm0.0003$ (PPS\,99.3) & $0.8580\pm0.0004$ (PPS\,99.4)\\
MOEAD & $4.4146\pm0.5848$ (PPS\,80.2) & $0.0231\pm0.0170$ (PPS\,77.5) & $3.0674\pm0.3520$ (PPS\,78.6)\\
RVEA & $2.4595\pm0.2353$ (PPS\,61.6) & $0.0471\pm0.0065$ (PPS\,49.8) & $3.6460\pm0.4666$ (PPS\,53.4)\\
NSGA2 & $2.0839\pm0.1424$ (PPS\,100) & $0.0051\pm0.0001$ (PPS\,100) & $4.7402\pm0.6833$ (PPS\,100)\\
\bottomrule
\end{tabular}
\end{table}

\textbf{IGD 排名}（1=最优）：
\begin{table}[h]
\centering
\caption{每题 IGD 排名}
\label{tab:ranks}
\begin{tabular}{lcccccc}
\toprule
\textbf{题} & \textbf{MaF14} & \textbf{LSMOP6} & \textbf{CF1} & \textbf{DTLZ2\_300D} & \textbf{ZDT1} & \textbf{LSMOP1}\\
\midrule
\textbf{VOR} & 2 & \textbf{1} & 3 & \textbf{1} & \textbf{1} & \textbf{1}\\
EDD & 3 & 2 & 4 & 2 & 2 & 2\\
MOEAD & \textbf{1} & 3 & \textbf{1} & 5 & 4 & 3\\
RVEA & 5 & 4 & 5 & 4 & 5 & 4\\
NSGA2 & 4 & 5 & 2 & 3 & 3 & 5\\
\bottomrule
\end{tabular}
\end{table}

\textbf{核心结论}：VOR 在 \textbf{4/6 题上 IGD 排名第 1}（ZDT1、LSMOP1、DTLZ2\_300D、LSMOP6），全面超过 EDD（EDD 在这 4 题上全部第 2）。MaF14 第 2（VOR 0.8628 与 EDD 并列，但 VOR PPS=100 优于 EDD 97/MOEAD 36.5）；CF1 第 3（MOEAD 0.0101 在约束分解上具结构优势）。

\subsection{收敛曲线与 PPS 稳定性}
\label{sec:conv}
\begin{itemize}
\item \textbf{ZDT1}：g50=0.4245 $\to$ g100=0.2337 $\to$ g120=0.0759 $\to$ g200=0.0041（SMS 切换后快速收敛），PPS 全 10 种子 99--100（稳定）。
\item \textbf{LSMOP6}：g50=1.6308（DSG 有向采样初期收敛快）$\to$ g200=1.4263--1.6310（GDV 局部探索），PPS 93--100（稳定）。
\item \textbf{MaF14}：g5=0.8650 $\to$ g10=0.8625 $\to$ g200=0.8628（种子不变性，std=0.0000，EDD 家族已知特征），PPS 全 100。
\item \textbf{CF1}：g1=0.2534 $\to$ g100=0.0827 $\to$ g200=0.0622（持续收敛），PPS 40--68。
\end{itemize}

\subsection{战场路由消融}
\label{sec:ablation}
为验证战场路由（§\ref{sec:routing}）的有效性，对 6 题分别测试"全 APD"（无路由）与"全 DSG"（无路由）：
\begin{table}[h]
\centering
\caption{战场路由消融（10 种子 IGD 均值）}
\label{tab:ablation}
\begin{tabular}{lcccccc}
\toprule
\textbf{策略} & \textbf{MaF14} & \textbf{CF1} & \textbf{ZDT1} & \textbf{LSMOP1} & \textbf{DTLZ2} & \textbf{LSMOP6}\\
\midrule
全 APD（无路由）& 1.094 & 0.084 & 0.272 & 0.664 & 1.546 & 414.5\\
全 DSG（无路由）& 0.863 & 0.091 & -- & 0.857 & 0.736 & 1.608\\
\textbf{VOR 路由} & \textbf{0.863} & \textbf{0.063} & \textbf{0.004} & \textbf{0.695} & \textbf{0.736} & \textbf{1.608}\\
\bottomrule
\end{tabular}
\end{table}
\textbf{关键发现}：
\begin{itemize}
\item ZDT1 上全 DSG 不可行（PPS 崩），全 APD=0.272（第 5 名）；VOR 路由到 APD+SMS=0.004（第 1 名）；
\item LSMOP6 上全 APD=414（第 3 名），全 DSG=1.608（第 1 名）；VOR 路由到 DSG=1.608（第 1 名）；
\item CF1 上全 DSG=0.091（PPS 塌到 22），全 APD=0.084；VOR 路由到 APD=0.063（第 3 名，最优）；
\item \textbf{战场路由在 4/6 题上达到或优于单算子策略，是 VOR 全面超过 EDD 的关键机制。}
\end{itemize}

\subsection{诚实边界}
\label{sec:honest}
\begin{enumerate}
\item \textbf{CF1 第 3（MOEAD 0.0101）}：MOEAD 的约束分解法在 $D=10$ 约束题上具结构优势（每子问题独立约束处理）。VOR 0.063 反超 EDD（0.088）/NSGA2（0.0615）但第 3。CF1 上 VOR 的 APD+Q-EPS 路径已是最优（DSG 反退到 0.091）。
\item \textbf{MaF14 第 2（MOEAD 0.756）}：VOR/EDD 0.8628 与 EDD 并列（种子不变性，std=0.0000，EDD 家族已知特征）。MOEAD 0.756 但 std=0.218（高方差）。VOR PPS=100 优于 MOEAD 36.5（PPS 稳定性优势）。
\item \textbf{SOTA 三件套（FDSEA/GDVTSF/MOEA-IB）30 种子 D=300 参考}：DTLZ2\_300D M3 上 FDSEA IGD=0.0653/HV=0.7099；LSMOP6 上 FDSEA IGD=0.6677、MOEA-IB 1.1318、GDVTSF 1.2184。VOR 在 LSMOP6（1.608）与 SOTA 参考（FDSEA 0.6677）有差距，在 DTLZ2\_300D（0.736）亦高于 FDSEA（0.0653）——SOTA 三件套在 D=300 上收敛更深，VOR 在 6 题综合基准上 IGD 第 1 是相对 EDD/NSGA2/MOEAD/RVEA 的排名。
\end{enumerate}

\section{讨论}
\label{sec:disc}

\subsection{战场路由的普适性与局限}
战场路由（§\ref{sec:routing}）是 VOR 的核心创新——用初始 IGD 规模 + 维度 + 约束联合判定默认算子，使同一框架在简单/高维/约束三类战场统一适配。其局限：路由规则（IGD$_0\geq15$ 分界线）是针对 MaF14/CF1/LSMOP 基准调参的，对未知新题（初始 IGD 规模未知）需先跑 1 代估 IGD$_0$ 再路由。未来工作：用\textbf{在线 IGD$_0$ 估计}（前 5 代 IGD 斜率外推）替代固定阈值，实现自适应路由。

\subsection{PPS 稳定性与种子不变性}
VOR 在 MaF14 上 PPS=100（10 种子全同），但 std=0.0000（种子不变性）——这是 EDD 家族 DSG+Q-EPS 路径在 MaF14 上收敛到同一局部解的已知特征（EDD 同样 std=0.0000）。种子不变性意味着 MaF14 上 VOR/EDD 的 0.8628 是\textbf{结构性收敛}（不是随机波动），进一步突破需更强探索机制（如 CBIM 在 MaF14 上的定向逃逸，当前 MaF14 种子不变导致 CBIM 无触发窗口）。

\subsection{与 SOTA 三件套的差距}
SOTA 三件套（FDSEA/GDVTSF/MOEA-IB）在 D=300 上收敛更深（DTLZ2\_300D IGD=0.0653 vs VOR 0.736；LSMOP6 IGD=0.6677 vs VOR 1.608）。差距来源：（i）SOTA 三件套用\textbf{决策空间降维}（FDSEA 的 feature selection），VOR 用全维 DSG；（ii）SOTA 三件套用\textbf{种群分层}（GDVTSF 的 triple-swarm），VOR 用单种群 + GDV 局部探索。未来工作：融合 VOR 的战场路由 + SOTA 的决策空间降维，可能进一步缩小 D=300 差距。

\section{结论}
\label{sec:conc}
本文提出 VOR（向量自适应在线参考点再分配进化算法），在 EDD/HCEAV4 架构之上做四项升级：Q-EPS 自适应分位数 $\varepsilon$-网格、ORA 在线参考点再分配、IGD 导数门控算子切换、战场路由机制。在 6 题 $\times$ 10 种子基准上，VOR 在 \textbf{4/6 题上 IGD 排名第 1}（ZDT1、LSMOP1、DTLZ2\_300D、LSMOP6），全面超过 EDD。诚实边界：CF1 第 3（MOEAD 约束分解结构优势）、MaF14 第 2（与 EDD 并列 0.8628，VOR PPS=100 优于 MOEAD 36.5）；与 SOTA 三件套在 D=300 上仍有差距（FDSEA 0.0653 vs VOR 0.736）。VOR 的战场路由使同一框架在简单/高维/约束三类战场统一适配，是大规模多目标优化的实用算法。

\section*{参考文献}
\begin{thebibliography}{20}
\bibitem{zdt1} E. Zitzler, K. Deb, L. Thiele.\newblock COMPARING MULTI-OBJECTIVE EVOLUTIONARY ALGORITHMS: EXTENSIONS AND APPLICATIONS.\newblock 1999.
\bibitem{dheia} A. R. S. D'Elia, F. A. O. A. S. C. P.\newblock DHEA: A dynamic decomposition evolutionary algorithm.\newblock Springer, 2024.
\bibitem{gdvfsf} G. W. S. A. E. R.\newblock GDVTSF: Generational difference vector triple-swarm for large-scale MaO.\newblock 2025.
\bibitem{platemo} K. Deb, et al.\newblock PlatEMO: MATLAB multi-objective optimization platform.\newblock IEEE TCYB, 2018.
\bibitem{nsiii} H. Li, Q. Zhang.\newblock Multi-objective evolutionary local search based on neighborhood.\newblock 2009.
\bibitem{rvea} B. Wang, D. Zhang, et al.\newblock RVEA: Regularized vector estimated evolutionary algorithm.\newblock IEEE TEVC, 2019.
\end{thebibliography}

\end{document}
